% Outline: Introduction (1 page cap; the ONLY section where bullets are allowed)
\section{Introduction}

When a coding agent talks to a language model, most of what travels
over the wire is not the user's question. Intercepting a production coding
agent (OpenCode) with an echo endpoint, we measured a single agent-loop
request. It declared 170 tools whose JSON schemas occupied 147\,kB, 67\% of
that payload, before the model read a word of the task. The share is a
range, not a constant. Within that agent's own tool set it runs 17--57\%
of request bytes (median 23.5\%), falling as conversation accumulates
around a schema block that is resent whole each turn; across every tool set
we captured the range widens to 18--66\%. Figures use
canonical serialization; raw wire bytes differ by roughly two points. No length predictor in the line we survey reads any of it.

Reading it would matter because schedulers want to know how long
each request will run. Shortest-job-first ordering needs an output-length
estimate, and a line of work has learned such estimates from prompt text
alone~\cite{fuLTR2024,tao2026ranking}. The traffic we captured violates the assumptions
behind both the scalar and the prompt-only designs. A third of its requests
carry no tools (title generation and other utility calls), each client's
schema is byte-constant across turns of the same call kind, and the tool
vocabulary in the wild is open-ended, so a learned component must handle
tools it has never seen. In our held-out
evaluation, 95.5\% of test requests carried a tool combination never seen in
training, and every tool-bearing request captured from the live agent was
composed entirely of tools foreign to the training corpus.

Our deployment context is an API gateway. Gateways are an established tier
in classical service infrastructure, where one layer owns load balancing,
admission control, caching and failure isolation, and operators now place
that same tier in front of LLM traffic with little modification. Agent
requests in this study reach the model through such a gateway (VeloxMesh),
which owns routing, admission and provider health, and which consults our
scoring service at admission time (Fig.~\ref{fig:arch}). The gateway treats
scoring as optional by design: a documented 15\,ms budget bounds the
decision, and on timeout or error the request proceeds in arrival order.
The scheduling work we survey optimizes inside the engine and does not
address this budget, so a learned predictor must fit inside it, recalibrate
it, or move off the synchronous path. This paper asks which of the gateway's
classical mechanisms remain useful for agentic LLM traffic, and whether a
learned signal can operate within the budget the gateway grants it.

We argue for two design changes. The request should be read as
\emph{text}: a ranker that consumes prompt and schema text orders unseen-tool
requests with little observed $\tau$ difference across the two adequately
powered unseen-tool strata (equivalence untested), and none of the identity
or scalar features we tested can rank tools they have never seen. Prediction
should also carry its own error bounds: our reliability gate assigns each
request a confidence measured per cold-start stratum and abstains,
preserving arrival order, wherever that measurement does not exist.

This paper contributes:
\begin{itemize}
  \item a workload characterization of live agent traffic: across captured
  tool sets, 18--66\% of
        request bytes are unread tool schema (67\% in the agent's 170-tool
        configuration), about 2\,kB per tool in its native sets,
        resent every turn; 33\% zero-tool requests; schemas byte-constant
        within a call kind;
        and a tool set entirely out of vocabulary relative to the corpus we
        train on;
  \item schema-text ranking: $+0.19$~$\tau$ over a grid-searched scalar
        baseline, with no drop detected across the two adequately powered
        unseen-tool strata (equivalence untested),
        while schema-identity
        baselines gain $+0.008$ (within noise) and \emph{lose} accuracy when
        forced to use the identity feature;
  \item a decision-budget coverage curve: neither un-truncating the schema
        nor caching its encodings makes the CPU scorer fit the documented
        15\,ms default, under which 34.6\% of decisions silently fail open;
        a timeout matched to the measured GPU scorer (50--75\,ms) restores
        98--100\% coverage, permitting that much decision latency on the
        critical path of a scored request rather than incurring it;
  \item a contention threshold for learned ordering: with the engine's
        running-batch cap at its default the queue never forms and no policy
        separates, but capping it at sixteen puts 57--64\% of requests behind
        another, and there the learned ranker beats arrival order by 4.2\%
        while ordering by prompt length runs 8.5\% slower than not reordering
        at all;
  \item a serving-level attribution: switching the scheduler implementation
        alone moved mean latency 44\% with arrival order held fixed, and
        prefix caching a further 17--23\%, in both contention regimes;
  \item an evidence-based abstaining gate: per-stratum confidences derived
        from measured ranking quality replace a hardcoded constant that
        overstates reliability by $+0.25$--$+0.46$ everywhere; the lowest
        observed ranking quality occurred on \emph{new combinations of
        familiar tools} ($n{=}78$, below our reporting threshold).
\end{itemize}

We evaluate on tool-calling corpora and on a live gateway stack
serving Qwen3.5-9B, with pre-registered criteria fixed before each
measurement. At serving time our pre-registered superiority test fails: no ordering
policy separates from another by more than 2\%, while swapping the scheduler
implementation alone cuts mean latency by 44\%. We report this attribution
result in place of the improvement we expected.

The remainder of this paper is organized as follows.
Section~II provides serving background and defines our vocabulary;
Section~III positions this work against prediction-based, prediction-free,
and imprecise-information scheduling; Section~IV details data, models, and
pre-registered criteria; Section~V presents workload characterization,
ranking, cold-start, cost, and gate results; Section~VI discusses design
consequences and limitations; Section~VII concludes.

